\chapter{Vertex Operator Algebras, Modular Sewing, and Logarithmic CFT}
\label{ch:voa-modular-sewing-logarithmic-cft}

\ConventionsEuclidean

\section{Chiral Algebra Data}

Two-dimensional conformal field theory has an additional algebraic structure:
local holomorphic fields may be multiplied by operator products that depend
only on one complex coordinate.  A vertex operator algebra is the precise
algebraic encoding of this chiral operator product.

\begin{definition}[Vertex operator algebra]
\label{def:vertex-operator-algebra}
A vertex operator algebra consists of the following data.
\begin{enumerate}
\item A complex vector space
\[
  V=\bigoplus_{n\in\mathbb Z}V_n
\]
with finite-dimensional graded pieces and \(V_n=0\) for \(n\) sufficiently
negative.
\item A vacuum vector \(\mathbf 1\in V_0\).
\item A state-field map
\[
  Y(\,\cdot\,,z):V\longrightarrow \operatorname{End}(V)[[z,z^{-1}]],
  \qquad
  a\longmapsto Y(a,z)=\sum_{m\in\mathbb Z}a_m z^{-m-1}.
\]
\item A conformal vector \(\omega\in V_2\), whose modes
\[
  Y(\omega,z)=\sum_{n\in\mathbb Z}L_n z^{-n-2}
\]
obey the Virasoro algebra
\[
  [L_m,L_n]=(m-n)L_{m+n}
  +\frac{c}{12}(m^3-m)\delta_{m+n,0}.
\]
\end{enumerate}
These data satisfy:
\[
  Y(\mathbf 1,z)=\operatorname{id}_V,\qquad
  Y(a,z)\mathbf 1\in a+zV[[z]],
\]
\[
  [L_{-1},Y(a,z)]=\partial_zY(a,z),
  \qquad
  L_0a=na\quad\text{for }a\in V_n,
\]
and locality:
for every \(a,b\in V\), there is \(N\ge0\) such that
\[
  (z-w)^N [Y(a,z),Y(b,w)]=0
\]
as a formal distribution on \(V\).
\end{definition}

The modes \(a_m\) are operators on \(V\).  If \(a\) has conformal weight
\(\Delta_a\), then \(a_m\) has degree \(\Delta_a-m-1\).  The vacuum axiom says
that \(a\) is recovered from the field \(Y(a,z)\) by acting on the vacuum and
then taking the value at \(z=0\):
\[
  a=\left. Y(a,z)\mathbf 1\right|_{z=0}.
\]
The locality axiom is the algebraic form of the chiral OPE.  It implies the
Jacobi identity for modes, and conversely the Jacobi identity packages
locality, associativity, and commutativity of OPE expansions.

\begin{proposition}[OPE from VOA locality]
\label{prop:voa-ope-from-locality}
For \(a,b\in V\), the product \(Y(a,z)Y(b,w)\) has an expansion in the region
\(|z|>|w|\) of the form
\[
  Y(a,z)Y(b,w)
  =
  \sum_{m\in\mathbb Z}Y(a_m b,w)(z-w)^{-m-1},
\]
where for each matrix element only finitely many positive powers of
\((z-w)^{-1}\) occur.
\end{proposition}

\begin{proof}
By lower truncation, \(a_m b=0\) for all sufficiently large \(m\).  Locality
gives \(N\) such that
\[
  (z-w)^N Y(a,z)Y(b,w)=(z-w)^N Y(b,w)Y(a,z).
\]
Taking the formal residue
\[
  a_m b=\operatorname{Res}_{z-w}(z-w)^mY(a,z)b
\]
defines the coefficient of \((z-w)^{-m-1}\).  Multiplying the proposed OPE by
\((z-w)^N\), both sides become the same formal Laurent series because the
residues of all principal parts agree and the remaining part is regular at
\(z=w\).  The truncation of \(a_m b\) gives the stated finiteness.
\end{proof}

\section{Modules, Characters, and Chiral Blocks}

A \(V\)-module \(M\) is a graded vector space with fields
\[
  Y_M(a,z)=\sum_{m\in\mathbb Z}a^M_m z^{-m-1}
\]
obeying the same vacuum, translation, and locality identities with one entry
acting on \(M\).  The conformal vector gives an operator \(L_0\) on \(M\).
For an ordinary positive-energy module,
\[
  M=\bigoplus_{\lambda\in\mathbb C}M_\lambda,\qquad
  \dim M_\lambda<\infty,
\]
and the set of \(\lambda\)'s is bounded below in every coset modulo
\(\mathbb Z\).  The character is
\begin{equation}
  \chi_M(\tau)
  =
  \operatorname{Tr}_M q^{L_0-c/24},
  \qquad
  q=\ee^{2\pi\ii\tau},
  \qquad
  \operatorname{Im}\tau>0 .
  \label{eq:voa-character}
\end{equation}
The shift by \(c/24\) is the cylinder vacuum energy fixed by the conformal map
from the plane to the circle.

Let \(C\) be a compact Riemann surface with marked points \(p_i\), local
coordinates \(z_i\) at the marked points, and \(V\)-modules \(M_i\).  A chiral
conformal block is a linear functional
\[
  \mathcal F:\bigotimes_i M_i\longrightarrow\mathbb C
\]
that obeys the Ward identities for meromorphic chiral fields on \(C\).  For
the stress tensor, the identity says that inserting a meromorphic vector field
regular away from the \(p_i\)'s acts by the sum of the corresponding Virasoro
modes at the marked points and gives zero.  For a general chiral state
\(a\in V\), the meromorphic one-form valued in the field \(a\) gives the same
condition with modes of \(Y_{M_i}(a,z_i)\).

On the sphere with three marked points, conformal blocks are intertwiners.
If \(M_i,M_j,M_k\) are simple modules, the fusion coefficient is
\[
  N_{ij}{}^k
  :=
  \dim \operatorname{CB}_{0,3}(M_i,M_j,M_k^\vee),
\]
where \(M_k^\vee\) is the contragredient module.  This definition is intrinsic:
it counts chiral three-point structures, not a particular coordinate
presentation of a three-point function.

\section{Sewing and Higher-Genus Blocks}

Sewing is the operation that constructs higher-genus surfaces and higher-point
blocks from lower-genus pieces.  Let \(C_1\) and \(C_2\) have marked points
with coordinates \(z\) and \(w\).  Removing disks and imposing
\[
  zw=q,\qquad |q|<1,
\]
glues the two punctures into an annulus.  If the two glued punctures carry a
module \(M_a\) and its dual \(M_a^\vee\), a sewn chiral block has the form
\begin{equation}
  \mathcal F_{\rm sewn}(q)
  =
  \sum_{\alpha}
  q^{h_a+n_\alpha-c/24}
  \mathcal F_1(\cdots,e_\alpha)
  \mathcal F_2(e^\alpha,\cdots).
  \label{eq:chiral-block-sewing}
\end{equation}
Here \(\{e_\alpha\}\) is a homogeneous basis of \(M_a\), \(e^\alpha\) is the
dual basis with respect to the invariant pairing, and
\(L_0 e_\alpha=(h_a+n_\alpha)e_\alpha\).  The exponent records propagation
along the sewn annulus.

\begin{definition}[Sewing-compatible rational chiral theory]
\label{def:sewing-compatible-rational-chiral-theory}
For this chapter, a rational chiral theory is a VOA \(V\) with finitely many
simple ordinary modules \(\{M_i\}_{i\in I}\), semisimple module category,
finite-dimensional conformal-block spaces, nondegenerate two-point pairings,
and convergent sewing expansions \eqref{eq:chiral-block-sewing} in a
neighborhood of \(q=0\).  The analytic continuations of these blocks along the
moduli space of pointed curves give a projectively flat vector bundle of
conformal blocks.
\end{definition}

The mapping class group acts on this bundle.  In genus zero the elementary
moves are associativity and braiding of OPE channels.  In genus one the
generators are
\[
  S:\tau\mapsto -1/\tau,\qquad
  T:\tau\mapsto\tau+1.
\]
The compatibility of sewing with these moves is the Moore--Seiberg system:
pentagon, hexagon, and modular identities for associativity, braiding, and
twist data.  The modern algebraic home of this system is a modular tensor
category when the module category is semisimple.

\section{Modular Tensor Category and Verlinde Formula}

Let \(\mathcal C\) be the semisimple ribbon category of modules of a rational
chiral theory.  Its simple objects are \(M_i\), with tensor product
\[
  M_i\boxtimes M_j\cong \bigoplus_k N_{ij}{}^k M_k .
\]
The matrices \((N_i)_j{}^k=N_{ij}{}^k\) are the fusion matrices.  The category
has a braiding, duals, and twists \(\theta_i=\ee^{2\pi\ii h_i}\).  Its
modular \(S\)-matrix is the Hopf-link pairing normalized so that the vacuum
label \(0\) corresponds to \(V\).  In a unitary theory \(S\) is unitary and
\[
  d_i:=\frac{S_{i0}}{S_{00}}
\]
is the quantum dimension of \(M_i\).

\begin{quotedtheorem}[Modularity of rational chiral characters]
\label{qthm:voa-character-modularity}
For a rational, \(C_2\)-cofinite, positive-energy VOA satisfying the analytic
finiteness hypotheses needed for sewing, the characters
\(\chi_i(\tau)\) of simple modules span a finite-dimensional representation
of \(SL(2,\mathbb Z)\):
\[
  \chi_i(-1/\tau)=\sum_j S_{ij}\chi_j(\tau),
  \qquad
  \chi_i(\tau+1)=\sum_j T_{ij}\chi_j(\tau).
\]
The representation agrees with the modular data of the semisimple module
category.
\end{quotedtheorem}

\begin{theorem}[Verlinde formula from modular diagonalization]
\label{thm:verlinde-from-modular-diagonalization}
Assume the fusion matrices are simultaneously diagonalized by a unitary
matrix \(S\) with \(S_{0m}\ne0\):
\[
  (N_i S)_{jm}
  =
  \frac{S_{im}}{S_{0m}}S_{jm}.
  \label{eq:fusion-diagonalization-by-s}
\]
Then
\begin{equation}
  N_{ij}{}^k
  =
  \sum_m
  \frac{S_{im}S_{jm}\overline{S_{km}}}{S_{0m}} .
  \label{eq:verlinde-formula}
\end{equation}
\end{theorem}

\begin{proof}
The coefficient \(N_{ij}{}^k\) is the \(j,k\) matrix element of \(N_i\):
\[
  N_{ij}{}^k=(N_i)_j{}^k.
\]
Equation~\eqref{eq:fusion-diagonalization-by-s} says
\[
  N_i
  =
  S D_i S^{-1},
  \qquad
  (D_i)_{mm}=\frac{S_{im}}{S_{0m}}.
\]
Since \(S\) is unitary, \(S^{-1}_{mk}=\overline{S_{km}}\).  Therefore
\[
  (N_i)_j{}^k
  =
  \sum_m S_{jm}\frac{S_{im}}{S_{0m}}\overline{S_{km}},
\]
which is \eqref{eq:verlinde-formula}.
\end{proof}

The theorem separates the deep input from the algebra.  The deep input is
that the same \(S\)-matrix obtained by the modular transformation of torus
characters diagonalizes the fusion ring.  Once this is granted, the formula is
linear algebra.

\section{The Ising Chiral Theory}

The Ising minimal model has central charge
\[
  c=\frac12
\]
and three simple modules:
\[
  \mathbf 1,\qquad \sigma,\qquad \varepsilon ,
\]
with conformal weights
\[
  h_{\mathbf 1}=0,\qquad h_\sigma=\frac1{16},\qquad
  h_\varepsilon=\frac12 .
\]
In the ordered basis \((\mathbf 1,\sigma,\varepsilon)\), the modular
\(S\)-matrix is
\begin{equation}
  S_{\rm Ising}
  =
  \frac12
  \begin{pmatrix}
    1 & \sqrt2 & 1\\
    \sqrt2 & 0 & -\sqrt2\\
    1 & -\sqrt2 & 1
  \end{pmatrix}.
  \label{eq:ising-modular-s}
\end{equation}
It is real, unitary, and obeys \(S^2=1\), because every simple object is
self-dual.  The quantum dimensions are
\[
  d_{\mathbf 1}=1,\qquad d_\sigma=\sqrt2,\qquad d_\varepsilon=1,
  \qquad
  \mathcal D^2=\sum_i d_i^2=4.
\]
Applying Theorem~\ref{thm:verlinde-from-modular-diagonalization} gives the
fusion rules
\[
  \varepsilon\boxtimes\varepsilon=\mathbf 1,\qquad
  \varepsilon\boxtimes\sigma=\sigma,\qquad
  \sigma\boxtimes\sigma=\mathbf 1\oplus\varepsilon .
\]

\begin{proof}
The identity \(S^2=1\) is a direct matrix multiplication using
\((\sqrt2)^2=2\).  The vacuum row gives the quantum dimensions by
\(d_i=S_{i0}/S_{00}\).  For example,
\[
  N_{\sigma\sigma}{}^{\mathbf 1}
  =
  \sum_m\frac{S_{\sigma m}S_{\sigma m}S_{\mathbf 1m}}{S_{\mathbf 1m}}
  =
  S_{\sigma\mathbf 1}^2+S_{\sigma\sigma}^2+S_{\sigma\varepsilon}^2
  =
  \frac12+0+\frac12=1,
\]
and
\[
  N_{\sigma\sigma}{}^\sigma
  =
  \sum_m\frac{S_{\sigma m}S_{\sigma m}S_{\sigma m}}{S_{\mathbf 1m}}
  =
  \frac{\sqrt2}{2}-\frac{\sqrt2}{2}=0.
\]
The same substitution gives
\(N_{\sigma\sigma}{}^\varepsilon=1\),
\(N_{\varepsilon\varepsilon}{}^{\mathbf 1}=1\), and
\(N_{\varepsilon\sigma}{}^\sigma=1\), with all other coefficients fixed by
commutativity and the unit object.
\end{proof}

\section{Full CFT and the Modular Bootstrap}

A full local CFT is not a chiral theory alone.  In a rational diagonal case,
one constructs a torus partition function from left and right characters:
\begin{equation}
  Z(\tau,\bar\tau)
  =
  \sum_{i,j}M_{ij}\chi_i(\tau)\overline{\chi_j(\tau)} .
  \label{eq:rational-torus-partition-function}
\end{equation}
The matrix \(M\) must have nonnegative integer entries, \(M_{00}=1\), and must
commute with the modular representation:
\[
  MS=SM,\qquad MT=TM .
  \label{eq:modular-invariant-matrix}
\]
These are genus-one sewing constraints.  Higher-genus sewing imposes further
compatibility between \(M\), chiral blocks, duality moves, and the pairing of
left and right chiral sectors.

For a unitary compact two-dimensional CFT with discrete spectrum on the
circle, the rectangular torus partition function
\[
  Z(\beta)=\operatorname{Tr}_{\mathcal H(S^1)}
  \exp\{-\beta(H-c_{\rm tot}/24)\},
  \qquad c_{\rm tot}=c+\bar c,
\]
obeys modular invariance
\[
  Z(\beta)=Z(4\pi^2/\beta)
\]
when the spatial circle has circumference \(2\pi\).  If the low-temperature
trace is dominated by a unique vacuum, then as \(\beta\to0^+\),
\begin{equation}
  \log Z(\beta)
  =
  \frac{\pi^2 c_{\rm tot}}{6\beta}+o(\beta^{-1}).
  \label{eq:cardy-partition-leading}
\end{equation}
The derivation is only the modular transformation:
low temperature at \(4\pi^2/\beta\) gives
\[
  Z(4\pi^2/\beta)
  =
  \exp\left\{\frac{4\pi^2}{\beta}\frac{c_{\rm tot}}{24}\right\}
  (1+o(1)),
\]
which is \eqref{eq:cardy-partition-leading}.  Tauberian hypotheses convert
this high-temperature partition-function statement into an asymptotic density
of states.  Those hypotheses are analytic input; modular invariance alone is
not a pointwise formula for individual degeneracies.

The modular bootstrap is the study of the constraints
\eqref{eq:modular-invariant-matrix}, reflection positivity, spin integrality,
and higher-genus sewing on the spectrum and OPE data.  In rational theories it
is finite-dimensional algebra.  In nonrational compact theories it is an
infinite-dimensional positivity and growth problem for torus and higher-genus
partition functions.

\section{Logarithmic Theories}

A logarithmic CFT is a conformal field theory in which the chiral module
category is not semisimple and \(L_0\) may have Jordan blocks.  If
\[
  L_0 v=hv+w,\qquad L_0w=hw,
\]
then cylinder evolution gives
\[
  q^{L_0-c/24}v
  =
  q^{h-c/24}(v+2\pi\ii\tau\,w).
\]
Correlation functions may therefore contain logarithms after mapping between
the plane and the cylinder.  Ordinary characters do not separate all modules;
one must include pseudo-traces and projective objects.  The category replacing
the semisimple modular tensor category is a finite tensor category with
braiding and twist, and the modular group acts on a larger space of
generalized characters and pseudo-trace functions.

The semisimple Verlinde formula
\eqref{eq:verlinde-formula} is then replaced by categorical statements about
the Grothendieck ring, projective modules, and modified traces.  The right
question is not whether the same matrix formula survives unchanged, but which
categorical pairing diagonalizes which quotient of the tensor product.  This
is one of the places where the algebraic formulation is essential: it names
the objects whose semisimplicity has failed and identifies the extra data
needed to recover modular sewing.

\section{Output}

The output of this chapter is a hierarchy of definitions and hypotheses:
The construction has the following order.  A VOA defines modules and chiral
blocks.  Chiral blocks sew over the moduli space of pointed curves.  In the
semisimple finite case, sewing data assemble into a modular tensor category.
The full CFT is then obtained by choosing a compatible pairing of left and
right chiral theories, equivalently by imposing the modular-invariance and
factorization constraints on torus and higher-genus amplitudes.
The Ising model gives the finite rational example.  Logarithmic theories show
which parts of the rational construction depend on semisimplicity and which
parts persist after replacing simple-module linear algebra by finite tensor
category data.
